\subsection{信道噪声}
\textbf{热噪声}\quad 高斯白噪声的一种. 其有效值与温度, 电阻性元件阻值和带宽有关:
\begin{equation}
    \overline{V}=\sqrt{4kTRB}.
\end{equation}

\textbf{白噪声}\quad 理想化噪声模型, 功率谱密度在整个频率范围均匀分布.
\begin{equation}
    P_n(\omega)=\frac{n_0}{2}.
\end{equation}
其自相关函数为
\begin{equation}
    R_n(\tau)=\frac{n_0}{2}\delta(\tau).
\end{equation}

\textbf{高斯白噪声}\quad 概率分布服从高斯分布的白噪声.

\textbf{带限白噪声}\quad 白噪声通过带宽有限信道或滤波器.

\textbf{低通白噪声}
\begin{gather}
    P_n(\omega)=\begin{cases}
        \frac{n_0}{2}, & \quad |f|\leq f_H;      \\
        0,             & \quad \text{otherwise}.
    \end{cases} \\
    R_n(\tau)=n_0f_H\frac{\sin2\pi f_H\tau}{2\pi f_H\tau}.
\end{gather}

\textbf{带通白噪声}
\begin{gather}
    P_n(\omega)=\begin{cases}
        \frac{n_0}{2}, & \quad f_c-\frac{B}{2}\leq|\omega|\leq f_c+\frac{B}{2}; \\
        0,             & \quad \text{otherwise}.
    \end{cases} \\
    R_n(\tau)=n_0B\frac{\sin\pi B\tau}{\pi B\tau}\cos2\pi f_c\tau.
\end{gather}
